\section*{Preface to Volume III}
\addcontentsline{toc}{section}{Preface to Volume III}

This volume deals with quantum field theories that are governed by
supersymmetry, a symmetry that unites particles of integer and half-integer
spin in common symmetry multiplets. These theories offer a possible way of
solving the `hierarchy problem,' the mystery of the enormous ratio of the
Planck mass to the 300 GeV energy scale of electroweak symmetry breaking.
Supersymmetry also has the quality of uniqueness that we search for in
fundamental physical theories. There is an infinite number of Lie groups that
can be used to combine particles of the same spin in ordinary symmetry
multiplets, but there are only eight kinds of supersymmetry in four spacetime
dimensions, of which only one, the simplest, could be directly relevant to
observed particles.

These are reasons enough to devote this third volume of \emph{The Quantum
Theory of Fields} to supersymmetry. In addition, the quantum field theories
based on supersymmetry have remarkable properties that are not found among
other field theories: some supersymmetric theories have couplings that are not
renormalized in any order of perturbation theory; other theories are finite;
and some even allow exact solutions. Indeed, much of the most interesting work
in quantum field theory over the past decade has been in the context of
supersymmetry.

Unfortunately, after a quarter century there is no direct evidence for
supersymmetry, as no pair of particles related by a supersymmetry
transformation has yet been discovered. There is just one significant piece of
indirect evidence for supersymmetry: the high-energy unification of the
$SU(3)$, $SU(2)$, and $U(1)$ gauge couplings works better with the extra
particles called for by supersymmetry than without them.

Nevertheless, because of the intrinsic attractiveness of supersymmetry and the
possibility it offers of resolving the hierarchy problem, I and many other
physicists are reasonably confident that supersymmetry will be found to be
relevant to the real world, and perhaps soon. Supersymmetry is a primary target
of experiments at high energy planned at existing accelerators, and at the
Large Hadron Collider under construction at the CERN laboratory.

After a historical introduction in Chapter 24, Chapters 25--27 present the
essential machinery of supersymmetric field theories: the structure of the
supersymmetry algebra and supersymmetry multiplets and the construction of
supersymmetric Lagrangians in general, and in particular for theories of chiral
and gauge superfields. Chapter 28 then uses this machinery to incorporate
supersymmetry in the standard model of electroweak and strong interactions, and
reviews experimental difficulties and opportunities. Chapters 29--32 deal with
topics that are mathematically more advanced: non-perturbative results,
supergraphs, supergravity, and supersymmetry in higher dimensions.

I have made the treatment of supersymmetry here as clear and self-contained as
I could. Wherever possible I take the reader through calculations, rather than
just reporting results from the literature. Where calculations are too lengthy
or complicated to be included in a book of this sort, especially in Chapter
28, I have tried to present simpler versions that give the reader an idea of
the physical issues involved.

I have made a point of including topics here that have generally not been
covered in earlier books, some because they are too new. These include: the use
of holomorphy to study perturbative and non-perturbative radiative corrections;
the calculation of central charges; gauge-mediated and anomaly-mediated
supersymmetry breaking; the Witten index; duality; the Seiberg--Witten
calculation of the effective Lagrangian in $N=2$ supersymmetric gauge theories;
supersymmetry breaking by modular fields; and a first look at the rapidly
developing topic of supersymmetry in higher dimensions, including theories
with $p$-branes.

On the other hand, I have shortened the usual treatment of two topics that
seemed to me to have been well covered in earlier books. One of these is the
use of supergraphs. Many of the previous applications of the supergraph
formalism in studying the general structure of radiative corrections can now
be handled more easily by using the arguments of holomorphy described in
Sections 27.6 and 29.3. The other is supergravity. In Sections 31.1--31.5 I have
given a detailed and self-contained treatment of supergravity in the weak-field
limit, which makes it clear why the ingredients of supergravity theories---the
graviton, gravitino, and auxiliary fields---are what they are, and which allows
us to derive some of the most important results of supergravity theory,
including the formula for the gravitino mass and for the gaugino masses
produced by anomaly-mediated supersymmetry breaking. In Section 31.6 I have
outlined the calculations that generalize supergravity theory to gravitational
fields of arbitrary strength, but these calculations are so lengthy and
unlovely that I was content to quote other sources for the results. However, in
Section 31.7 I have given a fuller than usual treatment of gravitationally
mediated supersymmetry breaking. I regret that I have not been able to include
exciting work of the past decade on supersymmetry related to string theory, but
string theory is beyond the scope of this book, and I did not want to report
results for which I had not provided a basis of explanation.

I have given citations both to the classic papers on supersymmetry and to
useful references on topics that are mentioned but not presented in detail in
this book. I did not always know who was responsible for material presented
here, and the mere absence of a citation should not be taken as a claim that
the material presented here is original, but some of it is. I hope that I have
improved on the original literature or standard textbook treatments in several
places, as for instance in the proof of the Coleman--Mandula theorem; in the
treatment of parity matrices in extended supersymmetry theories; in the
inclusion of new soft supersymmetry-breaking terms in the minimum
supersymmetric standard model; in the derivation of supercurrent sum rules; and
in the proof of the uniqueness of the Seiberg--Witten solution.

I have also supplied problems for each chapter. Some of these problems aim
simply at providing exercise in the use of techniques described in the
chapter; others are intended to suggest extensions of the results of the
chapter to a wider class of theories.

In teaching a course on supersymmetry, I have found that this book provides
enough material for a one-year course for graduate students. I intended that
this book should be accessible to students who are familiar with quantum field
theory at the level it is presented in the first two volumes of this treatise.
It is not assumed that the reader has gone through Volumes I and II, but for
the convenience of those fortunate readers who have done so I use the same
notation here, and give cross-references to material in Volumes I and II
wherever appropriate.

\begin{center}
  *\quad*\quad*
\end{center}

I must acknowledge my special intellectual debt to colleagues at the University
of Texas, notably Luis Boya, Phil Candelas, Bryce and Cecile De Witt, Willy
Fischler, Daniel Freed, Joaquim Gomis, Vadim Kaplunovsky, and especially Jacques
Distler. Also, Sally Dawson, Michael Dine, Michael Duff, Lawrence Hall, Hitoshi
Murayama, Joe Polchinski, Edward Witten, and Bruno Zumino gave valuable help
with special topics. Jonathan Evans read through the manuscript of this volume,
and made many valuable suggestions. For pointing out various errors in the
first printing of this book, I am greatly indebted to Stephen Adler, Jose
Espinora, Tony Gherghetta, and San Fu Tuan. For corrections to the first
printing of this volume I am indebted to several colleagues, especially Stephen
Adler. Thanks are due to Alyce Wilson, who prepared the illustrations, to Terry
Riley for finding countless books and articles, and to Jan Duffy for many
helps. I am grateful to Maureen Storey of Cambridge University Press for
working to ready this book for publication, and especially to my editor, Rufus
Neal, for his continued friendly good advice.

\vspace{2.5em}
\noindent
\begin{minipage}[t]{0.45\textwidth}
Austin, Texas\\
May, 1999
\end{minipage}
\hfill
\begin{minipage}[t]{0.45\textwidth}
\raggedleft
\textsc{Steven Weinberg}
\end{minipage}
